\documentclass[11pt]{amsart}
\usepackage[T1]{fontenc}
\usepackage{lmodern,microtype,amsmath,amssymb,amsthm,mathtools,enumitem}
\usepackage[margin=1in]{geometry}
\usepackage{booktabs}
\usepackage[hidelinks]{hyperref}
\newtheorem{mainthm}{Theorem}
\renewcommand{\themainthm}{\Alph{mainthm}}
\newtheorem{theorem}{Theorem}[section]
\newtheorem{proposition}[theorem]{Proposition}
\newtheorem{lemma}[theorem]{Lemma}
\newtheorem{corollary}[theorem]{Corollary}
\theoremstyle{definition}
\newtheorem{definition}[theorem]{Definition}
\newtheorem{question}[theorem]{Question}
\theoremstyle{remark}
\newtheorem{remark}[theorem]{Remark}
\newcommand{\bits}{\{-1,1\}}
\newcommand{\E}{\mathbb E}
\renewcommand{\P}{\mathbb P}
\newcommand{\R}{\mathbb R}
\newcommand{\HH}{\mathsf H}
\newcommand{\I}{\mathsf I}
\newcommand{\wh}[1]{\widehat{#1}}
\newcommand{\eps}{\varepsilon}
\DeclareMathOperator{\Stab}{Stab}
\DeclareMathOperator{\Ei}{E_1}
\setlist{itemsep=3pt,topsep=5pt}
\setlength{\emergencystretch}{1em}
\raggedbottom
\hypersetup{pdftitle={The Fourier entropy-influence inequality for symmetric Boolean functions},pdfauthor={Yitzchak Shmalo}}
\title[Fourier entropy of symmetric functions]{The Fourier entropy--influence inequality for symmetric Boolean functions}
\author{Yitzchak Shmalo}
\address{Einstein Institute of Mathematics, The Hebrew University of Jerusalem, Jerusalem, Israel}
\email{yitzchak.shmalo@gmail.com}
\date{9 November 2026}
\subjclass[2020]{42C10, 60C05, 68Q17, 06E30, 33C10}
\keywords{Fourier entropy--influence conjecture, symmetric Boolean function, noise stability, Krawtchouk polynomial, modified Bessel function}
\usepackage{xurl}
\begin{document}
\begin{abstract}
We prove that every symmetric Boolean function $f$ satisfies $\HH[f]\le c\,\I[f]$ with $c=931487/200000=4.657435$, improving the constant $12.04$ of O'Donnell, Wright and Zhou. We bound the entropy lost under a random coordinate restriction and average the bounds over the successive dimensions. The sharp one-step constant is $K=12-\frac92\log_23=4.8677\ldots$, attained by OR on three variables; retaining the dimension in the restriction argument gives a smaller entropy--influence constant. A noise-stability estimate for symmetric derivatives reduces the finite range to scalar inequalities. We certify them by outward-rounded interval arithmetic, use an exhaustive block-restriction census through dimension $20$, and close the infinite range with a decreasing analytic bound. A sharper estimate for the shift of the entropy kernel improves the finite certificates. Without computer verification, the same noise-stability estimate gives $\HH[f]\le6.191\,\I[f]$. AND and OR show that the optimal constant is at least $4$.
\end{abstract}
\maketitle

\section{Introduction}

For $f:\bits^n\to\bits$ write $f=\sum_{S\subseteq[n]}\wh f(S)\chi_S$ with $\chi_S(x)=\prod_{i\in S}x_i$. By Parseval the squared coefficients form a probability distribution. Its entropy and the total influence are
\[
 \HH[f]=\sum_{S}\wh f(S)^2\log_2\frac{1}{\wh f(S)^2},\qquad
 \I[f]=\sum_S|S|\,\wh f(S)^2=\sum_{i=1}^n\P_x\bigl[f(x)\ne f(x^{\oplus i})\bigr],
\]
where $x^{\oplus i}$ is $x$ with its $i$-th coordinate flipped. The Fourier entropy--influence conjecture of Friedgut and Kalai~\cite{FK96} asserts that $\HH[f]\le C\,\I[f]$ for an absolute constant $C$. It remains open. It is known for several classes of functions, among them symmetric functions~\cite{OWZ11,KMS12,CKLS16,ACKSW21,Shm26}.

A function is \emph{symmetric} if $f(x)$ depends only on the Hamming weight $|x|=\#\{i:x_i=-1\}$. O'Donnell, Wright and Zhou~\cite{OWZ11} proved the conjecture for symmetric functions with constant $12.04$.

\begin{mainthm}\label{thm:main}
Every symmetric Boolean function $f$ satisfies
\[
 \HH[f]\le c\,\I[f],\qquad c=\frac{931487}{200000}=4.657435<4.6575.
\]
\end{mainthm}

For a coordinate $i$ let $f_{i\to\pm1}$ be the two restrictions, $\I_i[f]$ the influence of $i$, and
\[
 \Delta_i[f]=\HH[f]-\tfrac12\bigl(\HH[f_{i\to1}]+\HH[f_{i\to-1}]\bigr).
\]
If $U_d$ bounds $\Delta_1/\I_1$ for symmetric functions in dimension $d$, then revealing coordinates one at a time gives $\HH[f]\le n^{-1}(\sum_{d=1}^nU_d)\I[f]$ (Lemma~\ref{lem:averaging}). A single uniform bound cannot be smaller than $K=12-\frac92\log_23$, because $\mathrm{OR}_3$ has $\Delta_1/\I_1=K$. Averaging over dimensions charges this obstruction only at dimension three. We obtain a stronger base by restricting several coordinates together through dimension $20$, then add the dimension-specific one-step bounds. This seeded average proves Theorem~\ref{thm:main}.

\begin{mainthm}\label{thm:restriction}
\begin{enumerate}[label=\textup{(\roman*)}]
\item For every $n\ge1$ and every symmetric Boolean $f$ on $n$ variables, $\Delta_1[f]\le K\,\I_1[f]$, with equality for $\mathrm{OR}_3$.
\item In high dimension the step is cheaper:
\[
 \limsup_{n\to\infty}\ \sup_{f}\ \frac{\Delta_1[f]}{\I_1[f]}\le C_*<4.637,
\]
where the supremum is over nonconstant symmetric $f$ on $n$ variables and
\[
 C_*=\frac{1}{\ln2}\Bigl[\frac{14}{5}+\frac{89}{50}\,e^{-10}I_0(10)+\frac{48}{23\sqrt{40\pi}}+\frac{e^{-z}}{z}\Bigr]=4.6364236659\ldots,\qquad z=40\Bigl(\frac{23}{48}\Bigr)^2.
\]
\end{enumerate}
\end{mainthm}

Part~(i) is computer-assisted. We check dimensions up to $15$ exhaustively. We verify the scalar inequalities separately in dimensions $16$ through $255$ and on $23$ blocks covering dimensions $256$ through $8{,}388{,}608$. The analytic estimate of Proposition~\ref{prop:large}, with parameter $a=1/30$, bounds every later one-step ratio by $4.654017$. Section~\ref{sec:computations} gives the certificates and the seeded ledger. Its largest upper bound is less than $4.657435$, at dimension $449$; this is a maximum of bounds, not a claimed value of $\HH/\I$.

Both parts rest on a noise-stability estimate for derivatives, which is of independent interest. Write $T_\lambda$ for the noise operator, $\Stab_\lambda[g]=\langle g,T_\lambda g\rangle=\sum_S\lambda^{|S|}\wh g(S)^2$, and $I_0$ for the modified Bessel function of order zero.

\begin{mainthm}\label{thm:stab}
Let $f$ be a symmetric Boolean function on $n\ge2$ variables, $N=n-1$, and let $g=D_nf$ be its derivative in the last coordinate, a function on $\bits^N$. For $0\le\lambda\le1$,
\[
 \Stab_\lambda[g]\le e^{-v}I_0(v)\,\|g\|_2^2,\qquad v=\frac{N(1-\lambda^2)}{4}.
\]
\end{mainthm}

Since $e^{-v}I_0(v)\sim(2\pi v)^{-1/2}$, this improves the derivative estimate used in~\cite{OWZ11} by a factor that tends to $2$ when $\lambda\to1$ and $v\to\infty$ together. The proof uses three facts: the derivative of a symmetric function is an alternating sum of indicators of Hamming slices; an alternating sum of selected values of a unimodal sequence is at most its largest value; and every atom of a sum of independent Bernoulli variables is bounded by Fourier inversion.

The restriction step loses most at the zero-level coefficient $\wh f(\emptyset)$, which appears in $\Delta_1$ but carries no influence. The asymptotic argument in part~(ii) bounds its contribution separately by $O(n^{-1/2})$; that is too crude in moderate dimensions. We instead bound it by a concave function of the mass $\mu_0$ that the derivative puts on level zero, and use a tangent of that function, so that the zero level enters the same linear program as the other levels (Lemma~\ref{lem:zero}). This is what makes the finite range tractable.

The noise-stability estimate also gives a bound that needs no computer at all.

\begin{mainthm}\label{thm:analytic}
Every symmetric Boolean function $f$ satisfies
\[
 \HH[f]\le\Bigl(\log_2(4+e)+2+\frac1{\ln2}\Bigr)\I[f]<6.191\,\I[f].
\]
\end{mainthm}

Here the entropy splits as $\HH=L+B$, where $L$ is the entropy of the level weights $W_k=\sum_{|S|=k}\wh f(S)^2$ and $B=\sum_kW_k\log_2\binom nk$ measures the spread inside the levels. The bound $L\le\log_2(4+e)\I$ holds for every Boolean function; it is Lemma~2.1 of~\cite{Shm26}. We prove $B<(2+1/\ln2)\I$ for nonconstant symmetric $f$ by integrating the envelope $e^{-v}I_0(v)$ against a harmonic-number kernel; a Frullani identity evaluates the main part exactly, and two inequalities for $e^{-x}I_0(x)$ and $e^{-x}I_1(x)$ control the rest.

The constant $4$ is a natural lower bound.

\begin{mainthm}\label{thm:lower}
For the OR function on $n$ variables (and for AND),
\[
 \frac{\HH[\mathrm{OR}_n]}{\I[\mathrm{OR}_n]}=4-\frac{4-2/\ln2}{n}+O(2^{-n}).
\]
Hence the optimal constant for symmetric functions lies in $[4,\,4.657435]$.
\end{mainthm}

That OR approaches $4$ is classical; we include the exact expression because AND and OR are the maximizers among all symmetric functions in every dimension we computed (Section~\ref{sec:computations}).

\subsection*{Notation}
Logarithms are natural unless written $\log_2$. Entropies are in bits. For a symmetric function we write $F(w)\in\bits$ for its value on weight $w$, $\wh f_k$ for the common coefficient at level $k$, and $W_k=\binom nk\wh f_k^{\,2}$. We put
\[
 \Phi(x)=e^{-x}I_0(x),\qquad G(x)=e^{-x}I_1(x)\qquad(x\ge0),
\]
so that $\Phi(x)=\frac1\pi\int_0^\pi e^{-x(1-\cos\theta)}\,d\theta$ and $\Phi'=G-\Phi$. Finally $\Ei(x)=\int_x^\infty e^{-u}u^{-1}\,du$ and $H_N=\sum_{i\le N}1/i$.

\section{The derivative of a symmetric function}\label{sec:split}

For $n\ge2$ let $N=n-1$ and $g=D_nf$, where $D_nf(x)=\tfrac12\bigl(f(x,1)-f(x,-1)\bigr)$ for $x\in\bits^N$. Then $\wh g(T)=\wh f(T\cup\{n\})$, so $g$ is symmetric with
\begin{equation}\label{eq:derivative}
 W_j[g]=\binom Nj\wh f_{j+1}^{\,2}=\frac{(j+1)W_{j+1}}{n},\qquad \|g\|_2^2=\sum_jW_j[g]=\I_n[f]=\frac{\I[f]}n.
\end{equation}
Its values are $g(x)=\tfrac12\bigl(F(w)-F(w+1)\bigr)$ with $w=|x|$. They are nonzero exactly on the weights $w_0<w_1<\dots$ where $F$ changes value, and there they alternate in sign. Up to a global sign,
\begin{equation}\label{eq:alternating}
 g(x)=\sum_{i}(-1)^i\,\mathbf 1[|x|=w_i].
\end{equation}
In particular $|g|=g^2$, so $\E|g|=\|g\|_2^2$. By symmetry the same holds for the derivative in any coordinate.

\section{Noise stability of derivatives}\label{sec:stab}

\begin{lemma}\label{lem:unimodal}
Let $(a_y)_{y\ge0}$ be a nonnegative unimodal sequence and $y_0<y_1<\dots<y_m$ integers. Then $\bigl|\sum_i(-1)^ia_{y_i}\bigr|\le\max_ya_y$.
\end{lemma}

\begin{proof}
Write $a_y=\int_0^\infty\mathbf 1[a_y>t]\,dt$. For each $t$ the set $\{y:a_y>t\}$ is an interval, so the selected indices in it form a block of consecutive $i$, whose alternating sum is $0$ or $\pm1$. The integrand vanishes for $t\ge\max a$. Integrate.
\end{proof}

\begin{lemma}\label{lem:atom}
Let $Z$ be a sum of $N$ independent Bernoulli variables with parameters in $\{p,1-p\}$, and $v=Np(1-p)$. Then $Z$ has a unimodal law and $\P[Z=y]\le\Phi(v)$ for every integer $y$.
\end{lemma}

\begin{proof}
Each Bernoulli law is log-concave and convolution preserves log-concavity of sequences supported on intervals, so the law of $Z$ is log-concave and in particular unimodal. Each summand has characteristic function of modulus $\bigl(1-4p(1-p)\sin^2(\theta/2)\bigr)^{1/2}$. Using $1-u\le e^{-u}$,
\[
 |\E e^{i\theta Z}|\le\exp\bigl(-2v\sin^2(\theta/2)\bigr)=e^{-v(1-\cos\theta)}.
\]
Fourier inversion gives $\P[Z=y]\le\frac1{2\pi}\int_{-\pi}^\pi|\E e^{i\theta Z}|\,d\theta\le\Phi(v)$.
\end{proof}

\begin{proof}[Proof of Theorem~\ref{thm:stab}]
Let $p=(1-\lambda)/2$. Given $x$ with $|x|=r$, the weight of a $\lambda$-correlated copy $y$ of $x$ is $Z_r=\mathrm{Bin}(r,1-p)+\mathrm{Bin}(N-r,p)$, of variance $Np(1-p)=v$. By~\eqref{eq:alternating},
\[
 \Stab_\lambda[g]=\E_x\Bigl[g(x)\sum_i(-1)^i\,\P\bigl[Z_{|x|}=w_i\bigr]\Bigr].
\]
By Lemmas~\ref{lem:unimodal} and~\ref{lem:atom} the inner sum has absolute value at most $\Phi(v)$. Hence $\Stab_\lambda[g]\le\Phi(v)\,\E|g|=\Phi(v)\|g\|_2^2$.
\end{proof}

\begin{remark}
For even $N$ the atom bound before the step $1-u\le e^{-u}$ is attained at $r=y=N/2$. The constant $\sup_v\sqrt v\,\Phi(v)=0.46882\ldots$ is the sharp uniform concentration constant for Bernoulli sums of Baillon, Cominetti and Vaisman~\cite{BCV16}; we use the full function $\Phi$ rather than this constant.
\end{remark}

\section{The restriction step}\label{sec:restriction}

\begin{lemma}\label{lem:induction}
Suppose $\Delta_1[f]\le K\,\I_1[f]$ for every symmetric $f$ on at most $m$ variables. Then $\HH[f]\le K\,\I[f]$ for every symmetric $f$ on at most $m$ variables.
\end{lemma}

\begin{proof}
Restrictions of a symmetric function are symmetric. For $i\ne1$ the influence of $i$ is the average of its influences in the two restrictions of $x_1$, so $\I[f]=\I_1[f]+\frac12(\I[f_{1\to1}]+\I[f_{1\to-1}])$. Induct on $n$, starting from constants.
\end{proof}

\begin{lemma}[Averaging over dimensions]\label{lem:averaging}
Suppose that $\Delta_1[g]\le U_d\I_1[g]$ for every symmetric Boolean function $g$ on $d$ variables. Then for every symmetric $f$ on $n\ge1$ variables,
\[
 \HH[f]\le\frac1n\sum_{d=1}^n U_d\,\I[f].
\]
More generally, if $\HH[g]\le C_m\I[g]$ in dimension $m$, then for $n\ge m$, $n\ge1$,
\begin{equation}\label{eq:seeded}
 \HH[f]\le\frac{mC_m+\sum_{d=m+1}^n U_d}{n}\,\I[f].
\end{equation}
Let $\Delta_{[k]}[f]$ be the entropy lost after fixing $k$ specified coordinates uniformly. If $\Delta_{[k]}[f]\le kB_{k,n}\I_1[f]$, one may choose the dimension bound
\begin{equation}\label{eq:blockrecurrence}
 C_n=\frac{kB_{k,n}+(n-k)C_{n-k}}{n}.
\end{equation}
One may minimize this choice over any available block sizes $k$.
\end{lemma}

\begin{proof}
A restriction of a symmetric weight table is a contiguous subtable, so it remains symmetric in the free coordinates. If $i$ is free and $Z$ fixes any set of other coordinates uniformly, then
\[
 \E_Z\I_i[f_Z]
 =\E_Z\P_{x_{\rm free}}[f_Z(x)\ne f_Z(x^{\oplus i})]
 =\P_x[f(x)\ne f(x^{\oplus i})]=\I_i[f].
\]
This is conditional expectation of the indicator of a changing edge. It includes constant restrictions.

Reveal the coordinates in a fixed order and let $F_d$ be the random function on the $d$ unrevealed coordinates. For the next coordinate $i_d$,
\[
 \HH[F_d]-\E[\HH[F_{d-1}]\mid F_d]\le U_d\I_{i_d}[F_d].
\]
On a branch with zero influence the symmetric function is constant, and both sides vanish. Taking expectations, influence preservation gives $\E\I_{i_d}[F_d]=\I_1[f]=\I[f]/n$. Summing over $d=1,\ldots,n$ telescopes to $\HH[f]$, since $F_0$ is constant. This proves the first assertion. For the second, stop at dimension $m$ and use $\E\HH[F_m]\le C_m\E\I[F_m]=mC_m\I[f]/n$. Finally, after fixing $k$ coordinates, $\E\I[f_Z]=(n-k)\I[f]/n$. Add the block entropy loss to $\E\HH[f_Z]\le C_{n-k}\E\I[f_Z]$ to obtain~\eqref{eq:blockrecurrence}. Constants, including dimension zero, have $\HH=\I=0$ and need no division. No sign assumption on the entropy losses is used.
\end{proof}

For symmetric $f$ on $n$ variables, the restrictions to $x_1=\pm1$ have coefficients $\wh f_s\pm\wh f_{s+1}$ on the $\binom Ns$ sets of size $s$ in the remaining coordinates, $N=n-1$. With $\varphi(u)=u\log_2(1/u)$ this gives the exact formulas
\begin{equation}\label{eq:delta}
 \Delta_1=\sum_{s=0}^{N}\binom Ns\,\mathcal D(\wh f_s,\wh f_{s+1}),\qquad \I_1=\sum_{s=0}^N\binom Ns\wh f_{s+1}^{\,2},
\end{equation}
\[
 \mathcal D(a,b)=\varphi(a^2)+\varphi(b^2)-\tfrac12\bigl[\varphi((a+b)^2)+\varphi((a-b)^2)\bigr].
\]
For $\mathrm{OR}_3$ one finds $\Delta_1/\I_1=12-\frac92\log_23$. From here on we work in natural units, $\mathcal D_e=(\ln2)\mathcal D$, and divide by $\ln2$ at the end. Throughout, $f$ is nonconstant, $g=D_1f$ and
\[
 \mu_s=\binom Ns\wh f_{s+1}^{\,2}/\I_1\qquad(0\le s\le N),
\]
a probability law on $\{0,\dots,N\}$ by~\eqref{eq:derivative}.

\subsection{A tangent bound}
Let $h(t)=\mathcal D_e(\sqrt t,1)$ for $t\ge0$. For $t>0$, $t\ne1$, direct differentiation gives
\[
 h''(t)=\frac{1}{2t^{3/2}}\ln\Bigl|\frac{\sqrt t-1}{\sqrt t+1}\Bigr|<0,
\]
$h'$ is continuous at $t=1$ and tends to $0$ from above, so $h$ is increasing and concave and $(x,y)\mapsto\mathcal D_e(\sqrt x,\sqrt y)=y\,h(x/y)$ is concave on the quadrant. Every tangent gives $\mathcal D_e(a,b)\le h'(t)a^2+[h(t)-th'(t)]b^2$. Taking the tangent at $t_s=(n-s)/s$ for the term $s\ge1$ of~\eqref{eq:delta}, and using $\binom Ns/\binom N{s-1}=(n-s)/s$ to move the $a^2$ part to the neighbouring index, one obtains
\begin{equation}\label{eq:tangent}
 \frac{(\ln 2)\Delta_1}{\I_1}\le L_0+\mu_0\kappa(\tfrac1n)+\sum_{s=1}^N\mu_s\Bigl[\Psi\bigl(\tfrac sn\bigr)+\kappa\bigl(\tfrac{s+1}n\bigr)-\kappa\bigl(\tfrac sn\bigr)\Bigr],\qquad L_0=\frac{\mathcal D_e(\wh f_0,\wh f_1)}{\I_1},
\end{equation}
where $\Psi(x)=h\bigl(\frac{1-x}x\bigr)$ and $\kappa(x)=\frac{1-x}xh'\bigl(\frac{1-x}x\bigr)$, extended continuously by $\kappa(0)=1$, $\kappa(1)=0$. Both $\Psi$ and $\kappa\bigl(\frac{s+1}n\bigr)-\kappa\bigl(\frac sn\bigr)$ are explicit, so the right side of~\eqref{eq:tangent} is a linear function of $\mu$ plus the zero-level term $L_0$.

We also use $h(t)\le3+\ln(t+\frac16)$. With $p(y)=y^2\ln(y^2)$, $p(0)=0$, the second-difference identity (valid since $p''$ is locally integrable) and Jensen's inequality for the triangular density $1-|u|$ on $[-1,1]$, which has mean $0$ and second moment $\tfrac16$, give
\begin{equation}\label{eq:hbound}
 h(t)=3+\int_{-1}^1(1-|u|)\ln\bigl((\sqrt t+u)^2\bigr)\,du\le3+\ln\bigl(t+\tfrac16\bigr).
\end{equation}

\subsection{Three facts about the derivative}
\begin{enumerate}[label=(\alph*)]
\item Let $m=|\wh f_1|=|\E g|$. Since $|g|=g^2$, $m\le\I_1$. By~\eqref{eq:alternating} and Lemma~\ref{lem:unimodal} applied to the binomial weights, $m\le b_N:=\max_j\binom Nj2^{-N}$. By Parseval for $f$, $\wh f_0^{\,2}+nm^2\le1$.
\item Consecutive values of $g$ on adjacent weights have nonpositive product, so $\|D_jg\|_2^2=\frac12\bigl(\|g\|_2^2-\E[g(x)g(x^{\oplus j})]\bigr)\ge\frac12\|g\|_2^2$ for each $j$, and summing over $j$ gives $\E_\mu s\ge N/2$.
\item By Theorem~\ref{thm:stab}, $\sum_s\mu_s\lambda^s\le\Phi\bigl(N(1-\lambda^2)/4\bigr)$ for $0\le\lambda\le1$. With $\lambda=e^{-t/n}$,
\begin{equation}\label{eq:noise}
 \E_\mu e^{-ts/n}\le\Phi\bigl(v_n(t)\bigr),\qquad v_n(t)=\frac{N(1-e^{-2t/n})}4\qquad(t>0).
\end{equation}
\end{enumerate}

\subsection{The zero level}
\begin{lemma}\label{lem:zero}
Let $C=n-\frac16$ and $J_C(z)=z\bigl[3+\ln(1-Cz^2)-2\ln z\bigr]$ for $0<z<C^{-1/2}$. Then $L_0\le J_C(\mu_0)$. The function $J_C$ is concave, so for every $0<u<C^{-1/2}$
\[
 L_0\le d_C(u)+j_C(u)\,\mu_0,\qquad d_C(u)=\frac{2u}{1-Cu^2},\quad j_C(u)=J_C'(u)=1-2\ln u+\ln(1-Cu^2)-\frac{2Cu^2}{1-Cu^2}.
\]
\end{lemma}

\begin{proof}
If $m=0$ then $L_0=0$ and $\mu_0=0$. Otherwise $\mu_0=m^2/\I_1\le m\le n^{-1/2}$ by~(a). Since $\mathcal D_e(a,b)=b^2h(a^2/b^2)$, $h$ increases, and $\wh f_0^{\,2}\le1-nm^2$, the bound~\eqref{eq:hbound} gives
\[
 L_0=\mu_0\,h\bigl(\wh f_0^{\,2}/m^2\bigr)\le\mu_0\bigl[3+\ln(m^{-2}-C)\bigr]\le\mu_0\bigl[3+\ln(\mu_0^{-2}-C)\bigr]=J_C(\mu_0),
\]
using $m\ge\mu_0$ in the last step. Differentiating twice,
\[
 J_C''(z)=-\frac2z-\frac{6Cz}{1-Cz^2}-\frac{4C^2z^3}{(1-Cz^2)^2}<0,
\]
and the tangent at $u$ has intercept $J_C(u)-uJ_C'(u)=d_C(u)$.
\end{proof}

\begin{lemma}[The shift estimate]\label{lem:shift}
The function $\kappa$ is increasing on $[0,\frac12]$ and $\kappa(\frac12)=2\ln2$. For $0<x<1$, $x\ne\frac12$,
\begin{equation}\label{eq:sharpderivative}
 \kappa'(x)\le B(x):=2+2\ln\frac1{|1-2x|}.
\end{equation}
Consequently, whenever $0\le x\le x+d\le1$,
\[
 \kappa(x+d)-\kappa(x)\le d\bigl(4+2\ln(1/d)\bigr)
 \le d\bigl(6+4\ln(1/d)\bigr).
\]
\end{lemma}

\begin{proof}
For $t\ge1$, expansion of the entropy pair gives
\[
 h(t)=\ln t+3-\sum_{j\ge1}\frac{t^{-j}}{j(j+1)(2j+1)}.
\]
For $0<x<\frac12$, put $r=x/(1-x)$. Then
\[
 \kappa(x)=1+\sum_{j\ge1}\frac{r^j}{(j+1)(2j+1)},
\]
which increases; at $r=1$ the sum is $2\ln2-1$. Termwise differentiation on compact subintervals, $j/[(j+1)(2j+1)]\le1/(2j)$ and $(1-x)^{-2}\le4$ give
\[
 \kappa'(x)\le\frac{2[-\ln(1-r)]}{r}
 \le2[1-\ln(1-r)]\le2-2\ln(1-2x).
\]
The middle inequality follows from $(1-r)[-\ln(1-r)]\le r$, by integrating $1/(1-t)\le1/(1-r)$ on $[0,r]$. The last uses $1-r=(1-2x)/(1-x)\ge1-2x$.

For $x>\frac12$, put $t=(1-x)/x<1$. The identity $h(t)=t h(1/t)$ gives
\[
 \kappa=t(2-\ln t)-\sum_{j\ge1}\frac{t^{j+1}}{j(2j+1)},\qquad
 \frac{d\kappa}{dt}=1-\ln t-\sum_{j\ge1}\frac{(j+1)t^j}{j(2j+1)}.
\]
Since $(j+1)/[j(2j+1)]=1/(2j)+1/[2j(2j+1)]$ and $\sum_{j\ge1}1/[2j(2j+1)]\le\frac12$, we have
\[
 \frac{d\kappa}{dt}\ge\tfrac12-\ln t+\tfrac12\ln(1-t)\ge\tfrac12\ln(1-t).
\]
Multiplying by $dt/dx=-x^{-2}$ reverses the inequality. Thus
\[
 \kappa'(x)\le-\frac{\ln(1-t)}{2x^2}
 \le-2\ln(1-t)\le-2\ln(2x-1),
\]
which is stronger than~\eqref{eq:sharpderivative}. The explicit entropy formula, or these series at their endpoints, shows that $\kappa$ is absolutely continuous, with integrable logarithmic singularities in its derivative. On an interval of length $d$ the integral of $\ln(1/|1-2y|)$ is largest when the interval is centred at $1/2$, where it equals $d(1-\ln d)$. Integrating~\eqref{eq:sharpderivative} proves the stated shift bounds.
\end{proof}

\subsection{The finite range}
\begin{proposition}\label{prop:dual}
Let $n\ge2$, $C=n-\frac16$, $0<u<C^{-1/2}$, and let $\alpha$, $\beta\ge0$, $w_j\ge0$, $t_j>0$ $(1\le j\le J)$ be real numbers with
\begin{align*}
 \kappa(\tfrac1n)+j_C(u)&\le\alpha+\textstyle\sum_jw_j,\\
 \Psi\bigl(\tfrac sn\bigr)+\kappa\bigl(\tfrac{s+1}n\bigr)-\kappa\bigl(\tfrac sn\bigr)&\le\alpha-\beta\tfrac sn+\textstyle\sum_jw_je^{-t_js/n}\qquad(1\le s\le N).
\end{align*}
Then every nonconstant symmetric $f$ on $n$ variables satisfies
\[
 \frac{(\ln2)\Delta_1}{\I_1}\le\alpha+d_C(u)-\beta\frac{N}{2n}+\sum_jw_j\,\Phi\bigl(v_n(t_j)\bigr).
\]
\end{proposition}

\begin{proof}
Insert Lemma~\ref{lem:zero} into~\eqref{eq:tangent}, bound each coefficient of $\mu_s$ by the hypotheses, and use $\E_\mu s\ge N/2$ from~(b) and~\eqref{eq:noise}.
\end{proof}

Proposition~\ref{prop:dual} also certifies whole ranges of dimensions. Fix $256\le L\le n\le U$ and use $C=L-\frac16$, which only enlarges $J_C$, and $\kappa(\frac1n)\le\kappa(\frac1L)$. We bound the entire concave function $J_C$ before taking its tangent; no monotonicity of the tangent slope or intercept separately is required.

The sharper derivative estimate gives a separate shift bound on each rational cell $[l,r]$. Put $h_L=1/L$ and
\[
 x_0=\min\{r,\max\{l,(1-h_L)/2\}\},\qquad
 M_{[l,r]}=\mathcal G(x_0+h_L-\tfrac12)-\mathcal G(x_0-\tfrac12),
\]
where $\mathcal G(z)=4z-2z\ln(2|z|)$ for $z\ne0$, and $\mathcal G(0)=0$. Indeed $\mathcal G'(y-1/2)=B(y)$. For $x=s/n\in[l,r]$, Lemma~\ref{lem:shift} gives
\[
 \kappa(x+1/n)-\kappa(x)\le\int_x^{x+1/n}B(y)\,dy
 \le\int_x^{x+h_L}B(y)\,dy\le M_{[l,r]}.
\]
In the second integral we extend $B(y)=2-2\ln|1-2y|$ as far as $1+h_L$ if needed. It is nonnegative there for $L\ge256$. No value of $\kappa$ outside $[0,1]$ is used. The integral over a fixed length $h_L$ increases until its starting point is $(1-h_L)/2$ and decreases afterwards: its derivative is $B(x+h_L)-B(x)$, whose sign follows by comparing the distances from $1/2$. This proves the formula for the maximum over a cell, including cells crossing $1/2$.

Both $\Psi$ and $x\mapsto\alpha-\beta x+\sum_jw_je^{-t_jx}$ decrease. Therefore the zero test at $1/L$ and the cell inequalities
\begin{equation}\label{eq:cell}
 \Psi(l)+M_{[l,r]}\le\alpha-\beta r+\sum_jw_je^{-t_jr}
\end{equation}
on a rational covering of $[1/U,1]$ prove all the hypotheses. On the objective side, $N/(2n)\ge(L-1)/(2L)$ and $v_n(t)$ increases in $n$. To see the latter, set $z=2t/n$ and differentiate:
\[
 4\frac{d}{dn}v_n(t)=1-(1+z)e^{-z}+\frac zn e^{-z}>0.
\]
Since $\Phi$ decreases, a certified block bound is
\begin{equation}\label{eq:blockobjective}
 U_{[L,U]}=\frac{\alpha+d_{L-1/6}(u)-\beta(L-1)/(2L)+\sum_jw_j\Phi(v_L(t_j))}{\ln2}.
\end{equation}

\subsection{High dimension}
For the analytic range we use the pointwise majorant
\begin{equation}\label{eq:majorant}
 \Psi(x)\le U(x):=\Ei(40x)+\frac{553}{100}-\frac{273}{50}x+\frac{89}{50}e^{-20x}\qquad(0<x\le1).
\end{equation}
It is verified by outward-rounded interval arithmetic on $11{,}900$ cells covering $[10^{-4},1]$, using that $\Psi$ and $U$ are decreasing (Appendix~\ref{app:cert}), with minimum certified slack above $0.02288$; for $0<x\le10^{-4}$ one uses $\Psi(x)\le3-\ln x$, $\Ei(40x)\ge-\gamma-\ln(40x)$ and $e^{-20x}\ge1-20x$. We also use, with $r_0=\frac{23}{48}$,
\begin{equation}\label{eq:Phibound}
 \Phi(v)\le\frac{24}{23\sqrt{2\pi v}}+e^{-2r_0^2v},\qquad
 \bar J:=\frac{48}{23\sqrt{40\pi}}+\frac{e^{-z}}z\ge\int_{40}^\infty\Phi(t/2)\frac{dt}t,\quad z=40r_0^2.
\end{equation}
These follow from $\Phi(v)=\frac2\pi\int_0^{\pi/2}e^{-2v\sin^2u}du$, $\sin u\ge2r_0u$ on $[0,\frac12]$ and $\sin u\ge r_0$ on $[\frac12,\frac\pi2]$, and $\Ei(z)\le e^{-z}/z$.

\begin{proposition}\label{prop:large}
Fix $0<a\le\frac15$. For $n\ge40/a$ put $N=n-1$, $b=\sqrt{2/(\pi N)}$, $c=\frac Nn\cdot\frac{1-e^{-2a}}{2a}$, $v_T=N(1-e^{-2a})/4$ and $A=-\frac12-\ln a+a$. Then every nonconstant symmetric $f$ on $n$ variables satisfies $(\ln2)\Delta_1/\I_1\le R_a(n)$, where
\begin{multline*}
 R_a(n)=b\bigl[3+\ln(1+b^{-2})+2\ln2\bigr]+\frac{6+4\ln n}{n}+\frac{14}5+\frac{273}{100n}+\bar J\\
 +\Phi(20c)\ln\frac1c+e^aA\Bigl[\frac{24}{23\sqrt{2\pi v_T}}+e^{-2r_0^2v_T}\Bigr]+\frac{89}{50}\Phi\bigl(v_n(20)\bigr).
\end{multline*}
Every term of $R_a$ is nonincreasing for $n\ge40/a$. In particular, $a=1/30$ is valid for $n\ge1200$, and $R_{1/30}(8{,}388{,}608)/\ln2<4.654016341712388$.
\end{proposition}

\begin{proof}
By~(a) and~\eqref{eq:hbound} with $\wh f_0^{\,2}\le1$, $L_0\le m[3+\ln(1+m^{-2})]$, and $u\mapsto u[3+\ln(1+u^{-2})]$ increases. Also $\mu_0\le m$, $\kappa(\frac1n)\le\kappa(\frac12)=2\ln2$, and $b_N\le b$ (for even $N$ write $b_N=\frac2\pi\int_0^{\pi/2}\cos^Nu\,du$ and use $\cos u\le e^{-u^2/2}$; for odd $N$, $b_N=b_{N+1}$). So the zero-level terms of~\eqref{eq:tangent} are at most the first term of $R_a$, and by Lemma~\ref{lem:shift} the shifts contribute at most $(6+4\ln n)/n$.

For $s\ge1$ apply~\eqref{eq:majorant} at $x=s/n$. The linear part gives at most $\frac{553}{100}-\frac{273}{50}\E_\mu\frac sn\le\frac{14}5+\frac{273}{100n}$ by~(b), and the last part at most $\frac{89}{50}\Phi(v_n(20))$ by~\eqref{eq:noise}. For the $\Ei$ part let $P(t)=\sum_{s\ge1}\mu_se^{-ts/n}$, so that $\sum_{s\ge1}\mu_s\Ei(40s/n)=\int_{40}^\infty P(t)\frac{dt}t$ by Tonelli. Put $T=an\ge40$. For $40\le t\le T$, since $(1-e^{-2y})/(2y)$ decreases, $v_n(t)\ge ct/2$, and~\eqref{eq:noise} gives $P(t)\le\Phi(ct/2)$; hence
\[
 \int_{40}^TP(t)\frac{dt}t\le\int_{40c}^\infty\Phi(y/2)\frac{dy}y\le\bar J+\Phi(20c)\ln\frac1c.
\]
For $t\ge T$ and $s\ge1$, $e^{-ts/n}\le e^{-(t-T)/n}e^{-Ts/n}$, so $P(t)\le e^{-(t-T)/n}\Phi(v_T)$ and
\[
 \int_T^\infty P(t)\frac{dt}t\le\Phi(v_T)\int_a^\infty e^{a-y}\frac{dy}y=e^a\Ei(a)\Phi(v_T)\le e^aA\,\Phi(v_T),
\]
using $\Ei(a)\le-\gamma-\ln a+a$ and $\gamma>H_6-\ln7>\frac12$ (since $H_6=49/20$ and $\ln7<39/20$). Bound $\Phi(v_T)$ by~\eqref{eq:Phibound}. Removing the zero level from $P$ is what keeps this integral finite.

For the monotonicity, the first term equals $\sqrt{2/\pi}\,[3+2\ln2+\ln(1+\pi N/2)]/\sqrt N$, whose derivative in $N$ is negative; $(6+4\ln n)/n$ decreases for $n\ge1$; $c$ increases in $n$ and stays in $(0,1)$; and $v_T$ and $v_n(20)$ increase in $n$ while $\Phi$ decreases. Here $A>0$ is fixed, and $c=(1-1/n)(1-e^{-2a})/(2a)$ increases within $(0,1)$. Thus $\Phi(20c)\ln(1/c)$ is a product of nonnegative decreasing factors. The analytic upper bound for $\Phi(v_T)$ in~\eqref{eq:Phibound} decreases as well. More explicitly, writing $z=\pi N/2$, twice the derivative of the zero-level term, multiplied by $N^{3/2}/\sqrt{2/\pi}$, is
\[
 \frac{2z}{1+z}-[3+2\ln2+\ln(1+z)]<0.
\]
The only lower threshold used in the integral split is $T=an\ge40$. The bound $\Ei(a)\le-\gamma-\ln a+a$ follows by integrating $(1-e^{-x})/x\le1$; hence the whole proof holds for every stated $a$. The value at $a=1/30$ and $n=8{,}388{,}608$ is certified in interval arithmetic.
\end{proof}

\subsection{Proofs of Theorems~\ref{thm:restriction} and~\ref{thm:main}}
\begin{proof}[Proof of Theorem~\ref{thm:restriction}(i)]
\emph{Dimensions $n\le15$.} The Fourier coefficients of a symmetric $f$ are integers divided by $2^n$, given by Krawtchouk transforms. For each of the $\sum_{n\le15}2^{n+1}=131{,}068$ symmetric functions we compute $\Delta_1$ and $\I_1$ from~\eqref{eq:delta}, with every logarithm enclosed between rational bounds from the series of $\operatorname{artanh}$ with an explicit remainder. Equality $\Delta_1=K\I_1$ holds for $n=3$ at OR, AND and their negations, where the value is an exact combination of $\log_22$ and $\log_23$; for every other function the upper enclosure of $\Delta_1/\I_1$ is below $K$.

\emph{Dimensions $16\le n\le255$.} For each $n$ we apply Proposition~\ref{prop:dual} with $J=2$ and decimal rational parameters, checking all $n$ hypotheses in outward-rounded interval arithmetic at $65$ digits. These are $32{,}520$ inequalities in total.

\emph{Dimensions $256\le n\le8{,}388{,}608$.} We use the block form of Proposition~\ref{prop:dual} and the cell inequalities~\eqref{eq:cell} on $23$ contiguous blocks, with $65{,}139$ rational covering cells. The full table appears in Section~\ref{sec:computations}.

\emph{Dimensions $n\ge8{,}388{,}608$.} Proposition~\ref{prop:large} with $a=1/30$.

The bounds obtained are listed in Section~\ref{sec:computations}; each is below $K$.
\end{proof}

\begin{proof}[Proof of Theorem~\ref{thm:main}]
For $n\le20$ we use the certified values $B_{k,n}$ for $2\le k\le\min(5,n)$ described in Section~\ref{sec:computations}. Starting from $C_0=C_1=0$, define $C_n$ as the minimum of the right sides of~\eqref{eq:blockrecurrence} over these $k$. This gives $C_n<4$ for $n\le20$ and $C_{20}<3.944265467369$.

For $21\le n\le D:=8{,}388{,}608$, use the one-step bounds $U_d$ of Propositions~\ref{prop:dual} and~\ref{prop:large}, and put
\[
 S_{20}=20C_{20},\qquad S_n=S_{n-1}+U_n.
\]
Lemma~\ref{lem:averaging} gives $\HH[f]/\I[f]\le S_n/n$. The exact rational ledger in Section~\ref{sec:computations} verifies
\[
 \max_{20\le n\le D}\frac{S_n}{n}
 =4.6574344717526121503843276538\ldots<c,
\]
with its maximum at $n=449$. This decimal denotes the recorded rational upper bound, not an extremal value of $\HH/\I$. For $d>D$, Proposition~\ref{prop:large} gives $U_d\le4.654016341712388<c$. If $S_n\le nc$ and $U_{n+1}\le c$, then $S_{n+1}\le(n+1)c$. This proves every remaining dimension. Constants have $\HH=\I=0$.
\end{proof}

\begin{proof}[Proof of Theorem~\ref{thm:restriction}(ii)]
Take any sequence $n\to\infty$ with nonconstant symmetric $f_n$, and pass to a subsequence along which the law of $s/n$ under $\mu$ converges weakly to a law $\nu$ on $[0,1]$. By~(b), $\E_\nu X\ge\frac12$. Letting $n\to\infty$ in~\eqref{eq:noise},
\begin{equation}\label{eq:laplace}
 \E_\nu e^{-tX}\le\Phi(t/2)\qquad(t>0),
\end{equation}
and in particular $\nu(\{0\})=0$. In~\eqref{eq:tangent} the zero-level terms are at most $b_N[3+\ln(1+b_N^{-2})+2\ln2]\to0$, as in the proof of Proposition~\ref{prop:large}, and the shifts at most $(6+4\ln n)/n\to0$. The function $\Psi$ is unbounded at $0$ but satisfies $\Psi(x)\le3+\ln(1/x)$, and~\eqref{eq:noise} with $t=n/m$ gives $\mu[s\le m]\le e\,\Phi\bigl(N(1-e^{-2/m})/4\bigr)\le C\sqrt{m/N}$ for an absolute $C$. Summing over dyadic ranges, the contribution of $1\le s\le\delta n$ is $O(\sqrt\delta(1+\ln(1/\delta)))$ uniformly in $n$. This uniform integrability and weak convergence give $\sum_{s\ge1}\mu_s\Psi(s/n)\to\int\Psi\,d\nu$. Hence $\limsup(\ln2)\Delta_1/\I_1\le\sup_\nu\int\Psi\,d\nu$, the supremum over laws on $(0,1]$ with $\E_\nu X\ge\frac12$ and~\eqref{eq:laplace}. Since $\Ei(40x)=\int_{40}^\infty e^{-tx}\frac{dt}t$, the majorant~\eqref{eq:majorant}, Tonelli,~\eqref{eq:laplace} and $\E_\nu X\ge\frac12$ give
\[
 \int\Psi\,d\nu\le\int_{40}^\infty\Phi(t/2)\frac{dt}t+\frac{553}{100}-\frac{273}{100}+\frac{89}{50}\Phi(10)\le\bar J+\frac{14}5+\frac{89}{50}\Phi(10).
\]
Dividing by $\ln2$ gives $C_*$, and a positive series for $I_0(10)$ certifies $\Phi(10)<0.127833337164$, so $C_*<4.637$.
\end{proof}

\section{A bound without computation}\label{sec:analytic}

For symmetric $f$ every set of size $k$ carries the same coefficient, so
\begin{equation}\label{eq:split}
 \HH[f]=\sum_{k=0}^n\binom nk\wh f_k^{\,2}\log_2\frac{1}{\wh f_k^{\,2}}
 =\underbrace{\sum_kW_k\log_2\frac1{W_k}}_{L}+\underbrace{\sum_kW_k\log_2\binom nk}_{B}.
\end{equation}
The level weights satisfy $\sum_kW_k=1$ and $\sum_kkW_k=\I[f]$.

\begin{lemma}[{\cite[Lemma~2.1]{Shm26}}]\label{lem:level}
Every Boolean function satisfies $L\le\log_2(4+e)\,\I[f]$.
\end{lemma}

\begin{proof}
We may assume $f$ is not constant. Let $\eps=1-W_0>0$ and $Q(k)=W_k/\eps$ for $k\ge1$. Put $a=\log_2(1+e/4)$, so that $(2^a-1)2^{-ak}$, $k\ge1$, is a probability law. Gibbs' inequality against this law gives $\HH(Q)\le a\,\E_Qk-\log_2(2^a-1)$. Since $L=h_2(\eps)+\eps\HH(Q)$, $\eps\E_Qk=\I$ and $h_2(\eps)\le\eps\log_2(1/\eps)+\eps\log_2e$, we get
\[
 L\le a\I+\eps\bigl[\log_2(1/\eps)+\log_2e-\log_2(e/4)\bigr]=a\I+\eps\bigl[\log_2(1/\eps)+2\bigr].
\]
Let $p=\min\{\P[f=1],\P[f=-1]\}$, so $\eps=4p(1-p)$. The edge-isoperimetric inequality~\cite[Chapter~2]{OD14} gives $\I\ge2p\log_2(1/p)$, and
\[
 2p\log_2\frac1p-\eps\Bigl(1+\tfrac12\log_2\frac1\eps\Bigr)=2p\bigl[p\log_2\tfrac1p+(1-p)\log_2(1-p)\bigr]\ge0,
\]
because the bracket is concave on $[0,\tfrac12]$ (its second derivative is $(2p-1)/(p(1-p)\ln2)$) and vanishes at both ends. Hence $\eps[\log_2(1/\eps)+2]\le2\I$ and $L\le(a+2)\I=\log_2(4+e)\I$.
\end{proof}

\begin{lemma}\label{lem:reduction}
Let $f$ be symmetric on $n\ge2$ variables with $\I[f]>0$, and $N=n-1$. Then
\[
 B\le\frac{1+C_N}{\ln2}\,\I[f],\qquad
 C_N=\int_0^1\frac{\Phi\bigl(Nt(2-t)/4\bigr)-(1-t)^N}{t}\,dt.
\]
\end{lemma}

\begin{proof}
Using $\binom nk\le(en/k)^k$,
\[
 B\ln2\le\sum_{k\ge1}kW_k\Bigl(1+\ln\frac nk\Bigr)=\I\Bigl(1+\E_q\ln\frac n{j+1}\Bigr),\qquad q_j=\frac{(j+1)W_{j+1}}{\I},\ 0\le j\le N.
\]
By~\eqref{eq:derivative}, $q_j=W_j[g]/\|g\|_2^2$. Comparing the sum with an integral, $\ln\frac{n}{j+1}\le H_N-H_j=\int_0^1\frac{\rho^j-\rho^N}{1-\rho}\,d\rho$. Therefore
\[
 \E_q\ln\frac n{j+1}\le\int_0^1\frac{s(\rho)-\rho^N}{1-\rho}\,d\rho,\qquad s(\rho)=\sum_jq_j\rho^j=\frac{\Stab_\rho[g]}{\|g\|_2^2}.
\]
By Theorem~\ref{thm:stab}, $s(\rho)\le\Phi\bigl(N(1-\rho^2)/4\bigr)$. Substituting $\rho=1-t$, so that $1-\rho^2=t(2-t)$, gives the claim. The integrand of $C_N$ is $O(1)$ near $t=0$.
\end{proof}

We need three facts about $\Phi$ and $G$.

\begin{lemma}\label{lem:bessel}
\begin{enumerate}[label=\textup{(\roman*)}]
\item $\sqrt w\,\Phi(w)>1/\sqrt{2\pi}$ for every $w\ge\tfrac12$.
\item $J(X):=\int_0^Xu\,\bigl(-\Phi'(u)\bigr)\,du=XG(X)$ for every $X\ge0$.
\item $G(X)<1/\sqrt{2\pi X}$ for every $X>0$.
\end{enumerate}
\end{lemma}

\begin{proof}
(i) Let $h(w)=\sqrt w\,\Phi(w)$. Substituting $I_0=e^w\Phi$ into the Bessel equation $I_0''+I_0'/w-I_0=0$ gives
\[
 h''+2h'+\frac{h}{4w^2}=0.
\]
At a stationary point $h''=-h/(4w^2)<0$, so $h$ has no local minimum in $(0,\infty)$. At the left end $h(\tfrac12)\ge e^{-1/2}/\sqrt2>1/\sqrt{2\pi}$, since $I_0\ge1$ and $\pi>e$. As $w\to\infty$, the substitution $\theta=y/\sqrt w$ and dominated convergence (with $1-\cos\theta\ge2\theta^2/\pi^2$ on $[0,\pi]$) give $h(w)\to\frac1\pi\int_0^\infty e^{-y^2/2}dy=1/\sqrt{2\pi}$. If $h(w_0)\le1/\sqrt{2\pi}$ for some $w_0>\tfrac12$, then $h$ restricted to some $[\tfrac12,W]$ with $W>w_0$ attains its minimum at an interior point: if $h(w_0)<1/\sqrt{2\pi}$ choose $W$ with $h(W)>h(w_0)$, and if equality holds take any $W>w_0$, since then $h\ge h(w_0)$ on $[\tfrac12,\infty)$ by the first case. This interior minimum is a local minimum, a contradiction.

(ii) From $I_0'=I_1$ and $I_1'=I_0-I_1/x$ we get $\Phi'=G-\Phi$ and $G'=\Phi-G-G/x$. Hence $(xG)'=x(\Phi-G)=-x\Phi'$, and $xG(x)\to0$ as $x\to0$.

(iii) Integrating $I_1(X)=\frac1\pi\int_0^\pi e^{X\cos\theta}\cos\theta\,d\theta$ by parts and substituting $u=1-\cos\theta$,
\begin{align*}
 G(X)&=\frac X\pi\int_0^\pi e^{-X(1-\cos\theta)}\sin^2\theta\,d\theta
 =\frac X\pi\int_0^2e^{-Xu}\sqrt{u(2-u)}\,du\\
 &<\frac{\sqrt2X}{\pi}\int_0^\infty e^{-Xu}\sqrt u\,du=\frac1{\sqrt{2\pi X}}.\qedhere
\end{align*}
\end{proof}

\begin{proposition}\label{prop:CN}
$C_N<\ln4$ for every integer $N\ge1$. More precisely $C_1\le1$, $C_2\le59/64$, and for $N\ge3$
\[
 C_N<\ln4-\frac{2-\sqrt2}{\sqrt{\pi N}}+\frac1{2N}<\ln4.
\]
\end{proposition}

\begin{proof}
Split $C_N=T_1+T_2+T_3$ with
\[
 T_1=\int_0^1\frac{\Phi(\frac{Nt(2-t)}4)-\Phi(\frac{Nt}2)}{t}dt,\quad
 T_2=\int_0^1\frac{\Phi(\frac{Nt}2)-e^{-Nt}}{t}dt,\quad
 T_3=\int_0^1\frac{e^{-Nt}-(1-t)^N}{t}dt.
\]

\emph{The term $T_2$.} Write $\Phi(x)=\E_\theta e^{-x(1-\cos\theta)}$ with $\theta$ uniform on $[0,\pi]$. For each $\theta$, Frullani's integral gives $\int_0^\infty\bigl(e^{-Nt(1-\cos\theta)/2}-e^{-Nt}\bigr)\frac{dt}t=\ln\frac2{1-\cos\theta}$; the integrand is nonnegative, so by Tonelli
\[
 \int_0^\infty\frac{\Phi(Nt/2)-e^{-Nt}}{t}dt=\E_\theta\ln\frac{1}{\sin^2(\theta/2)}=-\frac4\pi\int_0^{\pi/2}\ln\sin u\,du=\ln4.
\]
On $t\ge1$ we have $Nt/2\ge\tfrac12$, so Lemma~\ref{lem:bessel}(i) gives $\int_1^\infty\Phi(Nt/2)\frac{dt}t\ge\int_1^\infty\frac{dt}{\sqrt{\pi N}\,t^{3/2}}=\frac2{\sqrt{\pi N}}$. Since $\int_1^\infty e^{-Nt}\frac{dt}t=\Ei(N)$,
\[
 T_2\le\ln4-\frac2{\sqrt{\pi N}}+\Ei(N).
\]

\emph{The term $T_1$.} The function $\Phi$ is completely monotone, so $-\Phi'$ is positive and decreasing. With $a=Nt(2-t)/4$ and $b=Nt/2$ we have $a\ge Nt/4$ and $b-a=Nt^2/4$, so $\Phi(a)-\Phi(b)\le\frac{Nt^2}4\bigl(-\Phi'(\frac{Nt}4)\bigr)$. By Lemma~\ref{lem:bessel}(ii) and~(iii),
\[
 T_1\le\int_0^1\frac{Nt}4\Bigl(-\Phi'\bigl(\tfrac{Nt}4\bigr)\Bigr)dt=\frac4NJ\bigl(\tfrac N4\bigr)=G\bigl(\tfrac N4\bigr)<\sqrt{\frac{2}{\pi N}}.
\]

\emph{The term $T_3$.} Since $\int_0^1\frac{1-(1-t)^N}{t}dt=H_N$ and $\int_0^1\frac{1-e^{-Nt}}tdt=\ln N+\gamma+\Ei(N)$,
\[
 T_3=H_N-\ln N-\gamma-\Ei(N)\ \ (\ge0).
\]

Adding, the $\Ei(N)$ terms cancel and $H_N-\ln N-\gamma<1/(2N)$ gives
\[
 C_N<\ln4+\frac{\sqrt2}{\sqrt{\pi N}}-\frac{2}{\sqrt{\pi N}}+\frac1{2N}.
\]
The last two terms together are negative when $N>\pi/\bigl(4(2-\sqrt2)^2\bigr)$, and this threshold is less than $3$ because $\pi<4<12(2-\sqrt2)^2$, the last inequality following from $\sqrt2<17/12$.

For $N=1$, $\Phi\le1$ gives $C_1\le\int_0^1\frac{1-(1-t)}tdt=1$. For $N=2$, $e^{-z}\le1-z+z^2/2$ together with $\E(1-\cos\theta)=1$ and $\E(1-\cos\theta)^2=\tfrac32$ gives $\Phi(v)\le1-v+\tfrac34v^2$. With $v=(1-\rho^2)/2$,
\[
 C_2\le\int_0^1\frac{\frac{1-\rho^2}2+\frac3{16}(1-\rho^2)^2}{1-\rho}\,d\rho=\frac34+\frac3{16}\cdot\frac{11}{12}=\frac{59}{64}.
\]
Both are below $\ln4$.
\end{proof}

\begin{proof}[Proof of Theorem~\ref{thm:analytic}]
For $n=1$ every Boolean function is constant or $\pm x_1$ and $\HH=0$. For $n\ge2$ and $\I>0$, Lemma~\ref{lem:reduction} and Proposition~\ref{prop:CN} give $B<\frac{1+\ln4}{\ln2}\I=(2+\frac1{\ln2})\I$, and Lemma~\ref{lem:level} gives $L\le\log_2(4+e)\I$. Add them using~\eqref{eq:split}.
\end{proof}

\begin{remark}
The two parts are not tight together. For majority on $n=2m+1$ variables, $B/\I\to2/\ln2$, so the within-level constant cannot go below $2.885$; and $L/\I\le2.5176$ for every symmetric function with $n\le16$, attained by AND and OR at $n=6$ (Section~\ref{sec:computations}). Numerically $C_N$ increases to $\ln4$: $C_1=0.6638$, $C_2=0.8911$, $C_3=0.9991$, $C_{10^6}=1.3855$.
\end{remark}

The majority limit in the remark is proved as follows. Write $n=2m+1$. The derivative of majority is the indicator of the middle slice of the $2m$-cube, and its normalized spectral weights are
\[
 q_{2j}=\binom{2j}{j}\binom{2m-2j}{m-j}4^{-m},\qquad q_{2j+1}=0.
\]
Thus $j$ has the beta-binomial law with parameters $(m,\tfrac12,\tfrac12)$, and $j/m$ converges in law to the arcsine law on $(0,1)$. By~\eqref{eq:derivative}, $B/\I$ is the expectation of $\log_2\binom{2m+1}{2j+1}/(2j+1)$. Stirling's formula gives the interior limit $h(x)/(x\ln2)$ with $h(x)=-x\ln x-(1-x)\ln(1-x)$. The estimates $q_{2j}\le C/\sqrt{(j+1)(m-j+1)}$ and $\ln\binom nk/k\le1+\ln(n/k)$ bound the contribution of $j\le\delta m$ by $C\sqrt\delta(1+\ln(1/\delta))+O(\ln(m)/\sqrt m)$, and the other end contributes less, so the expectations converge. With $x=\sin^2u$,
\[
 \E\frac{h(X)}X=\frac2\pi\int_0^{\pi/2}\bigl[-\ln(\sin^2u)-\cot^2u\ln(\cos^2u)\bigr]\,du=2,
\]
by $\int_0^{\pi/2}\ln\sin u\,du=-\frac\pi2\ln2$ and $\int_0^{\pi/2}\cot^2u\ln(\cos^2u)\,du=\pi(\ln2-1)$. Hence $B/\I\to2/\ln2$.

\section{AND and OR}\label{sec:lower}

\begin{proof}[Proof of Theorem~\ref{thm:lower}]
Let $f$ take the value $-1$ at exactly one point; this is OR or AND up to the signs of the inputs and the output, which change neither $\HH$ nor $\I$. Then $\wh f(\emptyset)=1-2^{1-n}$ and $\wh f(S)^2=4^{1-n}$ for $S\ne\emptyset$. Hence, with $W_0=(1-2^{1-n})^2$,
\[
 \HH[f]=W_0\log_2\frac1{W_0}+(2^n-1)\,4^{1-n}\cdot2(n-1),\qquad \I[f]=n\,2^{1-n}.
\]
The second term divided by $\I$ is $(2-2^{1-n})\frac{2(n-1)}n=4-\frac4n+O(2^{-n})$, and $W_0\log_2(1/W_0)=\frac{2^{2-n}}{\ln2}\bigl(1+O(2^{-n})\bigr)$ contributes $\frac{2}{n\ln2}+O(2^{-n})$.
\end{proof}

At $n=1024$ the ratio is $3.998911513751\ldots$.

\section{Computations}\label{sec:computations}

All entropy ratios below use bits. The one-step bounds $U_d$ bound $\Delta_1/\I_1$. In dimensions $n\le15$ exact enumeration gives $U_d\le K$, with equality only in dimension $3$. In dimensions $16\le n\le255$, Proposition~\ref{prop:dual} gives $U_d\le4.848905$, with the largest certified bound at $n=25$. The new block certificates give the following upper bounds, rounded upwards:
\begin{center}
\begin{tabular}{@{}lrr@{}}
\toprule
dimensions & one-step bound & cells\\
\midrule
$256\le n\le 294$ & $4.691397764337803$ & $989$ \\
$295\le n\le 339$ & $4.680231363778156$ & $1{,}316$ \\
$340\le n\le 390$ & $4.669771627054690$ & $1{,}774$ \\
$391\le n\le 449$ & $4.660834262012443$ & $1{,}598$ \\
$450\le n\le 517$ & $4.653071439774263$ & $1{,}433$ \\
$518\le n\le 595$ & $4.646337217003664$ & $1{,}323$ \\
$596\le n\le 685$ & $4.640615237927098$ & $1{,}242$ \\
$686\le n\le 788$ & $4.635610258129181$ & $1{,}354$ \\
$789\le n\le 907$ & $4.631114929966981$ & $1{,}376$ \\
$908\le n\le 1{,}044$ & $4.627616624048243$ & $1{,}311$ \\
$1{,}045\le n\le 1{,}201$ & $4.625229653517468$ & $1{,}284$ \\
$1{,}202\le n\le 1{,}382$ & $4.622146309801676$ & $2{,}648$ \\
$1{,}383\le n\le 1{,}590$ & $4.618720507476198$ & $2{,}798$ \\
$1{,}591\le n\le 1{,}829$ & $4.615719288500242$ & $2{,}558$ \\
$1{,}830\le n\le 2{,}104$ & $4.613032839769942$ & $2{,}397$ \\
$2{,}105\le n\le 2{,}420$ & $4.610575912037690$ & $2{,}299$ \\
$2{,}421\le n\le 2{,}499$ & $4.607196989753158$ & $2{,}232$ \\
$2{,}500\le n\le 9{,}999$ & $4.624103643056020$ & $2{,}294$ \\
$10{,}000\le n\le 16{,}384$ & $4.597828415472139$ & $3{,}261$ \\
$16{,}385\le n\le 65{,}535$ & $4.599485777517825$ & $4{,}450$ \\
$65{,}536\le n\le 262{,}143$ & $4.592695956170183$ & $6{,}162$ \\
$262{,}144\le n\le 1{,}048{,}575$ & $4.589502673515857$ & $8{,}511$ \\
$1{,}048{,}576\le n\le 8{,}388{,}608$ & $4.588487126623142$ & $10{,}529$ \\
\bottomrule
\end{tabular}
\end{center}
The blocks are contiguous and have $65{,}139$ rational cells in total. At and above $D=8{,}388{,}608$, the analytic bound is $U_d\le4.654016341712388$. The multipliers and tangent points are exact decimal rationals. Each proposed block intercept is increased by $1/1000$ before certification. The optimizer supplies proposals only; its output is not a proof. We check the zero coefficient, the cell inequalities~\eqref{eq:cell}, exact coverage of $[1/U,1]$, and the objective~\eqref{eq:blockobjective} by outward intervals at $65$ decimal digits. For $\Phi(v)$ we use its positive Bessel series with a geometric remainder when $v\le100$, and the upper bound~\eqref{eq:Phibound} otherwise. A second implementation checks all $32{,}520$ finite coefficient inequalities and all cells at $90$ digits, using different entropy formulas and a different primitive for the shift integral.

\subsection*{The finite block base}
For $2\le k\le5$ and $k\le n\le20$, let $B_{k,n}$ be the largest value of $\Delta_{[k]}/(k\I_1)$ over nonconstant symmetric functions. We enumerate every weight table up to output complementation. If $A_s=2^n\wh f_s$ and $E_n=4^n\HH[f]$, then
\[
 E_n=2n4^n-2\sum_s\binom ns A_s^2\log_2|A_s|,
\]
with zero summands omitted. Fixing $k$ coordinates, $r$ of weight one, leaves the subtable $F(r),\ldots,F(r+n-k)$; write $E_{n-k,r}$ for its scaled entropy. Then
\[
 4^n\Delta_{[k]}=E_n-2^k\sum_{r=0}^k\binom kr E_{n-k,r},\qquad
 4^nk\I_1=k2^{n+1}\sum_{w:F(w)\ne F(w+1)}\binom{n-1}w.
\]
Integer Krawtchouk transforms give the $A_s$. Parseval and the independent edge formula for influence are checked at every table. Rational logarithm enclosures come from
\[
 \ln u=2\sum_{j=0}^{29}\frac{z^{2j+1}}{2j+1}+R,\qquad
 z=\frac{u-1}{u+1},\quad 1\le u<2,\quad
 0\le R\le\frac{9}{4\cdot61\cdot3^{61}}.
\]
Range reduction and outward integer rounding give fixed-point enclosures with denominator $2^{56}$. A second complete pass strictly excludes every nonmaximizer; exact prime-log expressions prove all surviving ties. A separate integer and $144$-bit logarithm audit checks every surviving maximum. The census covers $2{,}097{,}150$ tables up to complementation through dimension $20$, including constants, and supplies all $70$ pairs $(k,n)$. The recurrence~\eqref{eq:blockrecurrence} gives
\[
 C_{20}=\frac{92232807018024312555229592687247931457194510016451}
 {23384026197294446691258957323460528314494920687616}
 <3.944265467369.
\]
This is a certified upper bound. It is not identified with the exact maximum of $\HH/\I$ in dimension $20$.

\subsection*{The averaged ledger}
Seed the sum with $20C_{20}$ and add the $U_d$ for $21\le d\le D$. Exact rational addition gives
\[
 \begin{gathered}
 \max_{20\le n\le D}\frac{20C_{20}+\sum_{d=21}^nU_d}{n}\\
 =\frac{373080419884801803846596023166747790405659360503078208212549379}
 {80104276753121998324396300035352914529481654272000000000000000}.
 \end{gathered}
\]
This equals $4.6574344717526121503843276538\ldots$ and the maximum occurs at $n=449$. The claimed $c=931487/200000$ is an upward rounding. We also sum by segment lengths and recompute the binding value and cutoff interface with $85$-digit decimal arithmetic. On a segment with constant bound $U$, the ratio has the form $U+A/n$, so its maximum lies at an endpoint. This gives a further check of the finite maximum. At $D$ the ledger is below $4.588797$, and every later one-step bound is below $c$ by Proposition~\ref{prop:large}. The seed $20C_{20}$ is essential to the description of this computation: the displayed ledger is not the pure sum $\sum_{d=1}^nU_d$.

The certificate scripts, rational multipliers, cells and ledger are included with the data. The computation guide gives the command for each statement and distinguishes the exhaustive certificates from numerical checks.

Over all symmetric functions with $n\le16$ (exact integer Fourier coefficients, floating-point logarithms) the repository's check script finds:
\begin{center}
\begin{tabular}{@{}ll@{}}
\toprule
quantity & maximum over symmetric $f$, $n\le16$\\
\midrule
$\HH/\I$ & $3.93027$, at AND/OR with $n=16$ \\
$L/\I$ & $2.51755$, at AND/OR with $n=6$ \\
$B/\I$ & $1.80986$ \\
$\Delta_1/\I_1$ & $4.86767$, at $\mathrm{OR}_3$ \\
\bottomrule
\end{tabular}
\end{center}
In every dimension $2\le n\le16$ the maximum of $\HH/\I$ is attained by AND and OR. The script also checks the integral $C_N$ by adaptive quadrature for $N\le200$ and for $N$ up to $10^6$, Lemma~\ref{lem:bessel} on dense grids, Theorem~\ref{thm:stab} on random symmetric derivatives with $n=20,40,80$, the closed form of Theorem~\ref{thm:lower}, and the constants $K$ and $C_*$ at $40$ digits. These checks support but are not part of the proofs; the computer-assisted parts use the exact enumeration for $n\le15$, the block census through dimension $20$, the interval certificates for Propositions~\ref{prop:dual} and~\ref{prop:large}, the exact rational ledger, and the certificate for~\eqref{eq:majorant}.

\section{Questions}

\begin{question}
Is the optimal constant for symmetric functions equal to $4$?
\end{question}

A single uniform restriction bound cannot be smaller than $K$, because of $\mathrm{OR}_3$. Lemma~\ref{lem:averaging} shows that this does not obstruct a smaller entropy--influence constant: one can retain the dimension-specific cost throughout the restriction chain.

For general Boolean functions, $K$ does not bound the entropy loss at a suitable coordinate. If every nonconstant $f$ had a coordinate with $\Delta_i[f]\le C\,\I_i[f]$, induction on the number of variables would give $\HH[f]\le C\,\I[f]$ for all Boolean $f$, so Hod's lower bound $C>6.454784$ for the entropy--influence conjecture~\cite{Hod17} already excludes every $C\le6.454784$. The obstruction is also visible at a single coordinate in small dimension. In the $\{0,1\}$ convention, consider the majority of three disjoint AND pairs,
\[
 f=\operatorname{Maj}_3(x_1\wedge x_2,\ x_3\wedge x_4,\ x_5\wedge x_6).
\]
After converting the output to $\{-1,1\}$, every coordinate has
\[
 \frac{\Delta_i[f]}{\I_i[f]}=5.1948716233791414624\ldots>K.
\]
The OR of eight disjoint AND pairs on sixteen variables gives, at every coordinate,
\[
 \frac{\Delta_i[f]}{\I_i[f]}=6.0037460565590882514\ldots>6.
\]
Both functions are transitive on their coordinates. Exact integer Walsh transforms and rational logarithm enclosures certify the two inequalities in the data.

\begin{question}
Is there an absolute constant $C$ such that every nonconstant Boolean function has a coordinate $i$ with $\I_i[f]>0$ and $\Delta_i[f]\le C\,\I_i[f]$?
\end{question}

An inequality of this kind is the key element of the approach of Gonz\'alez, MacManus and Pereyra~\cite{GMP26}, who note that it would imply the conjecture by induction and verify the conjecture for several classes of functions through stopping trees. A positive answer would prove the entropy--influence conjecture with constant $C$, so $C>6.454784$ by~\cite{Hod17}.

\appendix
\section{The scalar certificate for the majorant}\label{app:cert}

The interval certificate for~\eqref{eq:majorant} uses two elementary lower bounds for $\Ei$. For $0<z\le1$,
\[
 \Ei(z)\ge-0.578-\ln z+\sum_{k=1}^{8}\frac{(-1)^{k+1}z^k}{k\,k!},
\]
since an even truncation of the alternating series is a lower bound and $\gamma<H_{10}-\ln\frac{21}2<0.578$. The last bound holds because $H_n-\ln(n+\frac12)$ decreases to $\gamma$: the integral of the convex function $1/x$ over $[n+\frac12,n+\frac32]$ exceeds its value $1/(n+1)$ at the midpoint. For $z>1$,
\[
 \Ei(z)\ge e^{-z}\,\frac{z^2+8z+11}{z^3+9z^2+18z+6}.
\]
Indeed $\Ei(z)=e^{-z}\,\E[(z+Y)^{-1}]$ for $Y$ exponential with mean one, and for every quadratic polynomial $P$ we have $(z+y)^{-1}\ge2P(y)-(z+y)P(y)^2$, the difference being $[1-(z+y)P(y)]^2/(z+y)$. In the orthonormal basis $1$, $1-y$, $1-2y+y^2/2$ the compression of multiplication by $z+y$ to quadratic polynomials is
\[
 M=\begin{pmatrix}z+1&-1&0\\-1&z+3&-2\\0&-2&z+5\end{pmatrix},
\]
and optimizing over $P$ gives $(M^{-1})_{00}$, which is the displayed rational function. The entries of $M$ follow from $\E Y^k=k!$.

Both $\Psi$ and $U$ decrease: $h$ increases, and $U'(x)=-e^{-40x}/x-5.46-35.6\,e^{-20x}<0$. So an outward interval lower bound for $U(r)-\Psi(l)$ proves~\eqref{eq:majorant} on a whole cell $[l,r]$. For $0<x\le10^{-4}$ the bounds in the text give
\[
 U(x)-\Psi(x)\ge5.53+1.78-0.578-\ln40-3-(5.46+20\cdot1.78)\,10^{-4}>0.
\]
The value $\Phi(10)$ is enclosed using $I_0(10)=\sum_{k\ge0}25^k/(k!)^2$, with the terms through $k=60$ summed and the rest bounded by a geometric series of ratio at most $25/62^2$. The resulting bound is $C_*<4.636423665924$.

\section*{Acknowledgements and funding}
The final manuscript was audited by Claude Opus 5.5 (Anthropic) and by Codex agents (OpenAI), which checked the proofs, recomputed the numerical constants in high-precision arithmetic, and reran the finite computations in the repository.

The research presented in this paper was supported by the European Research
Council (ERC) under the European Union's Horizon Europe research and
innovation programme (grant agreement No.~101041711), by the Simons
Foundation as part of the Collaboration on the Mathematical and Scientific
Foundations of Deep Learning, by Heights Labs, by Convex Nexus Capital, by
the Israel Science Foundation (grant number 2258/19), and by the Israel
Science Foundation (ISF Grant 4101/25). Computational resources were
provided in part by Google Cloud through its education credits program
(credit 530550103).

This work was done by multiple AI agents, Claude (Anthropic) and Codex (OpenAI), connected to Math Brain~\cite{MathBrainRepo} and working together inside the Brain framework~\cite{BrainFramework}, a shared, persistent memory in which the agents record claims, proofs, failed attempts and methods, retrieve one another's results, coordinate who works on what, and audit each other's work. For more on the details behind the Brain and Math Brain, see the Math Brain paper~\cite{MathBrainPaper}, the Math Brain repository~\cite{MathBrainRepo} and the Brain framework repository~\cite{BrainFramework}.

\section*{Declaration of competing interest}
The author holds equity in Heights Labs and in Convex Nexus Capital and
has a management role at Convex Nexus Capital. These entities are included
in the funding declaration above.

\begin{samepage}
\section*{Declaration of generative AI and AI-assisted technologies}
The paper comes out of Math Brain, an ongoing project on the use of AI systems for research mathematics, and its results were developed with the Brain~\cite{BrainFramework} and Math Brain~\cite{MathBrainRepo}. In November 2026, Claude Opus 5.5 (Anthropic) and a team of Codex agents running OpenAI's GPT-6.1, all connected to the Brain, retrieved the spectral-level entropy lemma and related methods from the author's earlier work through Math Brain, derived the noise-stability estimate and the split of the entropy into level and within-level parts, proved the Bessel-function inequalities and the asymptotic restriction estimate, found the bound for the zero-level coefficient that closes the restriction step in every dimension, derived the averaging argument over successive dimensions and the sharper shift estimate, wrote and ran the interval certificates and the exhaustive computation for small dimensions, and searched the literature. Separate agents checked the averaging argument, the sharper shift estimate and the changed analytic cutoff, and recomputed the certificates with independently written code. Claude Opus 5.5 coordinated the agents, assembled the proofs, prepared the manuscript, and audited it, including the proofs, the numerical constants and the finite computations. These checks are not peer review. The author directed the work and is responsible for the content of the paper and for decisions concerning its dissemination.
\par
\end{samepage}

\section*{Data availability}
\url{https://github.com/yspennstate/symmetric-fei-constant}

The paper is part of the Math Brain project, whose public repository is \url{https://github.com/yspennstate/math-brain-public}~\cite{MathBrainRepo}.

\begin{thebibliography}{99}
\bibitem{ACKSW21}
S.~Arunachalam, S.~Chakraborty, M.~Kouck\'y, N.~Saurabh and R.~de~Wolf,
\emph{Improved bounds on Fourier entropy and min-entropy},
ACM Trans. Comput. Theory \textbf{13} (2021), no.~4, Article~22, 40~pp.
\href{https://doi.org/10.1145/3470860}{doi:10.1145/3470860}.

\bibitem{BCV16}
J.-B.~Baillon, R.~Cominetti and J.~Vaisman,
\emph{A sharp uniform bound for the distribution of sums of Bernoulli trials},
Combin. Probab. Comput. \textbf{25} (2016), no.~3, 352--361.
\href{https://doi.org/10.1017/S0963548315000127}{doi:10.1017/S0963548315000127}.

\bibitem{CKLS16}
S.~Chakraborty, R.~Kulkarni, S.~V.~Lokam and N.~Saurabh,
\emph{Upper bounds on Fourier entropy},
Theoret. Comput. Sci. \textbf{654} (2016), 92--112.
\href{https://doi.org/10.1016/j.tcs.2016.05.006}{doi:10.1016/j.tcs.2016.05.006}.

\bibitem{FK96}
E.~Friedgut and G.~Kalai,
\emph{Every monotone graph property has a sharp threshold},
Proc. Amer. Math. Soc. \textbf{124} (1996), no.~10, 2993--3002.
\href{https://doi.org/10.1090/S0002-9939-96-03732-X}{doi:10.1090/S0002-9939-96-03732-X}.

\bibitem{GMP26}
M.~J.~Gonz\'alez, P.~MacManus and M.~C.~Pereyra,
\emph{Further evidence towards the Fourier entropy-influence conjecture},
preprint (2026), \href{https://arxiv.org/abs/2606.00246}{arXiv:2606.00246}.

\bibitem{Hod17}
R.~Hod,
\emph{Improved lower bounds for the Fourier entropy/influence conjecture via lexicographic functions},
preprint (2017), \href{https://arxiv.org/abs/1711.00762}{arXiv:1711.00762}.

\bibitem{KMS12}
N.~Keller, E.~Mossel and T.~Schlank,
\emph{A note on the entropy/influence conjecture},
Discrete Math. \textbf{312} (2012), no.~22, 3364--3372.
\href{https://doi.org/10.1016/j.disc.2012.07.031}{doi:10.1016/j.disc.2012.07.031}.

\bibitem{OD14}
R.~O'Donnell,
\emph{Analysis of Boolean Functions},
Cambridge University Press, Cambridge, 2014.

\bibitem{OWZ11}
R.~O'Donnell, J.~Wright and Y.~Zhou,
\emph{The Fourier entropy--influence conjecture for certain classes of Boolean functions},
in: ICALP 2011, Lecture Notes in Comput. Sci. \textbf{6755}, Springer, Heidelberg, 2011, 330--341.
\href{https://doi.org/10.1007/978-3-642-22006-7_28}{doi:10.1007/978-3-642-22006-7\_28}.

\bibitem{Shm26}
Y.~Shmalo,
\emph{Fourier entropy of threshold functions and their compositions},
preprint, 2026.
\href{https://arxiv.org/abs/2610.16069}{arXiv:2610.16069}.

\bibitem{MathBrainRepo} Y.~Shmalo, \emph{Math Brain: public snapshot, code and tests}, 2026, \url{https://github.com/yspennstate/math-brain-public}.

\bibitem{MathBrainPaper} Y.~Shmalo, \emph{Math Brain: shared, auditable memory for multi-agent mathematical research}, October 2026, \url{https://github.com/yspennstate/math-brain-public/tree/main/papers/math-brain-2026-10-24}.

\bibitem{BrainFramework} Y.~Shmalo, \emph{Brain Framework}, public software repository, 2026, \url{https://github.com/yspennstate/brain-framework-public}.
\end{thebibliography}
\end{document}
